\section{From Observed Correlations to Genetic Parameters}
\label{sec:measurement-error}\label{sec:assortment-comparison}
\subsection{What Classical Measurement Error Assumes}
Let the standardized recorded outcome be
\[
 Y_i=\sqrt{\theta}P_i+\sqrt{1-\theta}U_i,
 \qquad 0<\theta\leq1.
\]
Here $\theta$ is reliability, not the noise share. If recording errors are
uncorrelated with all relevant relatives' true components and with one
another, and reliability is common, then
\[
 \operatorname{Corr}(Y_i,Y_j)=\theta\operatorname{Corr}(P_i,P_j)
 \quad(i\ne j).
\]
Attenuation occurs at the two observed endpoints, not at each intervening
generation. The genetic variance share of $Y$ is $\theta h^2$, whereas
$h^2$ refers to $P$. Appendix~\ref{app:measurement-error} derives these
relations and the different normalization required for parental averages.

These assumptions go beyond imperfect measurement. Shared classification
errors add covariance; unequal reliability produces the pair-specific
factor $\sqrt{\theta_i\theta_j}$. More fundamentally, the difference between
education and an asserted latent social ability can contain family resources,
institutional access, and preferences, rather than independent recording
noise. A plausible value of the fitted attenuation factor does not validate
these exclusions.

\subsection{Identification Depends on the Mating Model}
For the nonspousal kinships in Clark's regression, the phenotypic-assortment
form is
\[
 q_{n,D}=\theta h^2\left(\frac{1+m}{2}\right)^n
 \left(\frac{1+r}{1+m}\right)^D,
\]
where $n$ indexes genealogical distance and $D=1$ for the lineal rows.
The three coefficients of its logarithm identify $\theta h^2$, $m$, and
$r$, when the data supply the necessary distance and lineal variation.
Under phenotypic assortment, the restriction $m=h^2r$ then separates
$h^2=m/r$ from $\theta$. With known attenuation, that same restriction
instead reduces the number of freely fitted parameters. Arbitrary
regression coefficients need not yield admissible structural values.

Under genetic assortment, the lineal adjustment disappears. These
kinship correlations identify $m$ and the product $\theta h^2$, but not
reliability and heritability separately. Even the observed spouse
correlation is then $\theta h^2m$. The intuition is that untransmitted
true environmental variation and recording noise both dilute the genetic
signal without changing mate selection on genetic value. They cannot be
separated by naming one of them measurement error.

The book's equation~A3 \citep[Technical Appendix]{Clark2026} therefore requires a precise interpretation. Three
regression coefficients can be legitimate under phenotypic assortment
with unknown attenuation, provided recovery invokes the restriction and
checks admissibility. Their number alone is not an error. But the
intercept identifies $\theta h^2$, not latent heritability, and changing
the mating mechanism changes the permissible recovery. Illustrative fits
in Appendix~\ref{app:assortment-comparison} show how similar fit can accompany
different identification claims. They are descriptive comparisons, not
formal tests using a joint sampling covariance matrix.

\subsection{Binary Outcomes Are Not Continuous Phenotypes}
\label{sec:binary-outcomes}

Several of the outcomes to which Clark applies the continuous-trait kinship
equations are indicators: company-director status, higher education, and
literacy. The Pearson correlation of two indicators is a valid descriptive
statistic. The problem is the next step. Clark treats those correlations as
if they were generated by the same linear, Gaussian-style moment equations
that apply to a continuously varying phenotype, with one common attenuation
factor and the same restriction $m=h^2r$. He does not specify a
threshold/liability model, convert the observed binary correlations to latent
correlations, or show that a single attenuation factor is valid across
kinship types.

A calibrated Gaussian-threshold example shows the practical importance.
The latent generating parameters are $h^2=.60$, $r=.80$, $m=.48$, and
continuous reliability $\theta=.80$. We threshold the recorded continuous
outcome at the approximately 2.5\% educational prevalence in Clark's
historical cohorts. Appendix~\ref{app:binary-calibration} reports the design,
prevalence construction, and complete calibration table.

The continuous version recovers the generating parameters: in the education
calibration, the free-$\theta$ estimates are approximately $h^2=.60$,
$r=.80$, and $m=.48$. After thresholding at the observed prevalence, the
same estimator gives $h^2=.41$, $r=.91$, and $m=.37$ for 1780--1859 and
$h^2=.39$, $r=1.00$, and $m=.39$ for 1860--1919; the fitted log-correlation
RMSE remains only about .07. Thresholding changes the
mapping from latent to observed correlation in a way that varies by kinship
type, so a good-looking fit can coexist with a materially wrong structural
interpretation.

This problem concerns the observation equation even if the latent
inheritance model is correct. A threshold model can relate observed
prevalences and joint probabilities to latent correlations. One freely
estimated common attenuation factor cannot generally substitute for that
mapping; a separate free factor for every kinship would instead surrender
the cross-kinship restrictions being tested.

These simulations are not a formal rejection of the empirical model: they
use a specified latent Gaussian threshold process, independent pairs, and
prevalence calibration from unique individuals rather than the exact
kinship-pair samples. They establish the narrower point needed here. Clark
has not supplied the additional binary-outcome model or validation required
to interpret his indicator correlations as evidence for the continuous-trait genetic parameters he reports \citep[Table~A1 and equations~A1--A3]{Clark2026}.
\subsection{A Diagnostic in Clark's Reported Correlations}

The restriction derived in Section~\ref{sec:fisher-model} provides a concrete
way to examine Clark's own numbers. Table~\ref{tab:table45} reproduces
the raw sibling and parent-child (``Full Sibling'' and ``Child'') correlations
from the book's Table 4.5, for the two outcomes with two separate birth-year windows: occupational status and higher-education attainment \citep[p.~54]{Clark2026}.

\begin{table}[h]
\centering
\caption{Raw correlations from Clark's Table 4.5}
\label{tab:table45}
\begin{tabular}{lcccc}
\hline
 & \multicolumn{2}{c}{Occupational status} & \multicolumn{2}{c}{Higher education} \\
 & 1780--1859 & 1860--1919 & 1780--1859 & 1860--1919 \\
\hline
Full Sibling ($\hat\rho_S$)      & 0.584 & 0.496 & 0.558 & 0.359 \\
Child ($\hat\rho_{PO}$)          & 0.585 & 0.513 & 0.500 & 0.305 \\
\hline
$\hat\rho_{PO}-\hat\rho_S$       & $+0.001$ & $+0.017$ & $-0.058$ & $-0.054$ \\
\hline
\end{tabular}
\end{table}

If a common, positive, phenotype-specific attenuation factor scales
both correlations for a given outcome and period, it preserves the sign
of the parent--child minus sibling difference. Equation~\eqref{eq:overview-gap}
requires that population difference to be positive when assortment is
positive and $0<h^2<1$. The higher-education point estimates have the
opposite sign in both periods. Their binary scale and potentially different
kinship-sample prevalences require the additional observation model just
discussed; this table alone is not a test of a latent-liability restriction. This is a useful specification diagnostic,
but point estimates alone do not establish rejection: a formal test must
account for their joint sampling uncertainty and verify comparable
measurement and sampling across relationship types. If attenuation differs
between the two kinships, the simple sign comparison need not be valid.
The relevant exercise is therefore a joint test of the constrained
predictions, not an inference from overall fit alone.
